\chapter{Network requirements}
\label{ch:network}

\section{Ports}

\begin{longtable}{@{}p{0.12\textwidth}p{0.14\textwidth}p{0.62\textwidth}@{}}
\toprule
\textbf{Port} & \textbf{Protocol} & \textbf{Use}\\
\midrule
\endhead
54320 & UDP & Discovery: answers the plugin's \ui{Find}. Optional.\\
54321 & UDP & Audio: requests and replies of all sessions, and telemetry.\\
54322 & TCP & Control: the plugin lists, parameters, state, and \cmd{fxmenet-top}.\\
\bottomrule
\end{longtable}

All three can be changed with the server's options (chapter~\ref{ch:server}).
The client opens connections to the server; the server only needs these three
ports open to the client.

\section{Bandwidth}

Audio travels as 32-bit floating-point samples, uncompressed, in both
directions: about 2.8\,Mbit/s for a stereo track at 44.1\,kHz, 3.1\,Mbit/s at
48\,kHz, and the more for more channels. In the first stress tests, 24
stereo tracks used about 71\,Mbit/s each way and 32 tracks used 95\,Mbit/s
each way. A gigabit connection carries several hundred tracks of this kind
in theory; in practice, the server's CPU and the quality of the network
limit the number of tracks first.

Only the channels that the routing uses are sent
(chapter~\ref{ch:routing}).

\section{Security}

The traffic is not encrypted or authenticated. A token in each session keeps
stray packets out, but anyone on the network can talk to the server. Use
\fxme\ on a network you trust, and keep the ports closed on the internet.
